\subsection{The complete models for e = 14, 16, and 18}
\label{sec:u3-higher2}

\begin{proposition}[K17-HECKE-HIGHER2]\label{prop:hecke-higher2}
Retain the compact-unitary and classical full-space models of
Theorem~\ref{thm:u3-general} and Proposition~\ref{prop:gl2-general}.
For \(e=14,16,18\), let \(P_{e,p}\) denote the characteristic polynomial
of \(p^{e-1}T_p^+\) on the entire level-\(K\) compact-unitary coefficient
space of highest weight \((e-1,0,1-e)\), and let \(Q_{e,p}\) denote that
of the classical \(T_p\) on
\(S_{2e+1}(\Gamma_0(3),\chi_{-3})\).
The full dimensions and the complete invariant counts are
\[
\begin{array}{c|c|r|r|r|c|r}
e&k=e-1&\dim V_k^\Gamma&\dim V_{k-1}^\Gamma&
\dim\mathcal H_k^\Gamma&\text{swap }(+,-)&\dim S_{2e+1}\\ \hline
14&13&64&50&14&(8,6)&8\\
16&15&104&83&21&(12,9)&10\\
18&17&156&129&27&(15,12)&11.
\end{array}
\tag{HH.1}
\]
Put
\[
 E_{e,p}=1+p^e+p^{2e},\qquad
 N_{e,p}=P_{e,p}(E_{e,p})Q_{e,p}(1+p^{2e}).
\]
The exact common divisors of the full comparison values are
\[
\begin{aligned}
\gcd(N_{14,7},N_{14,13})
 &=2^{106}3^{101}5^7\,53^3\,419\,15083^3\,16519\\
 &\qquad{}\cdot7547263^3\,42800473^3,\\
\gcd(N_{16,7},N_{16,13})
 &=2^{123}3^{147}5^4\,11^4\,23\,401\,13687\\
 &\qquad{}\cdot36293^3\,38177^3\,25078104767246023^3,\\
\gcd(N_{18,7},N_{18,13})
 &=2^{177}3^{180}5^{17}\,7\,13^8\,19^2\,37^2\,73\\
 &\qquad{}\cdot3343^3\,7691\,8609\,18541^3
                 \,151871^3\,8361790201^3.
\end{aligned}
\tag{HH.2}
\]
Every displayed factor is prime. The only factors at least \(2^{31}\)
are \(25078104767246023\) at \(e=16\) and \(8361790201\) at \(e=18\).
The full six unitary and six classical matrices and characteristic
polynomials are recorded in the frozen exact data and reconstructed
by the accompanying standard-library certificates.
\end{proposition}

\begin{proof}
The class-number-one reduction and representation argument in
Theorem~\ref{thm:u3-general} are uniform in \(k\).
The complete invariant monomials satisfy
\[
 |\alpha|=|\beta|=k,\qquad
 \alpha_i-\beta_i\equiv0\pmod6.
\]
Their simultaneous permutation orbit sums therefore include exponent
differences \(0,\pm6,\pm12\) in the present degrees.
The exact enumeration gives the first two dimension columns in (HH.1).
The independent replay reconstructs the entire contraction matrix,
checks that its rank is its full target dimension, and checks
independence and annihilation for every saved harmonic basis column.
This proves the complete harmonic dimensions and the listed swap
eigenspace dimensions, without omitting any additional invariants.

The same \(57\) and \(183\) unitary neighbors are used.
For every saved harmonic basis vector \(h_j\), the scaled action is
still exactly
\[
 p^{e-1}\rho(U/\pi)h_j(z,w)=h_j(U^*z,U^tw).
\]
Swap still conjugates the two minuscule directions, and multiplication
of the swap-minus basis vectors by \(1+2\omega\) gives the same rational
form of the full matrix. The characteristic polynomial is integral
because the scaled operator preserves the full lattice of invariant
harmonic polynomials over \(\mathbf Z[\omega]\). No new direction,
normalization, or lattice hypothesis is introduced.

For \(k=15,17\), the discovery program chooses evaluation
points by rank screening modulo \(1000003\), then verifies the exact
nonzero determinant over \(\mathbf Q\). This is only a faster method
of finding a square identifying matrix: the unchanged independent
verifier reconstructs all evaluations over the integers and checks
the exact determinant. It then recomputes every transformed value
and verifies the two full matrix equations \(AM_0=B_0\) and
\(AM_1=B_1\). Thus it identifies the entire Hecke operator on the
already proved full space. Every saved operator is also checked at
two further independent point pairs.

The full-space classical basis is
\[
 Da^{\,2e-5-3j}b^j,\qquad
 0\le j\le\left\lfloor(2e-5)/3\right\rfloor ,
\]
by the uniform statement of Proposition~\ref{prop:gl2-general}.
It gives dimensions \(8,10,11\) and an integral unitriangular matrix
of initial Fourier coefficients. The new classical certificate uses
precision \(260\), greater than the maximum identifying index
\(13\cdot11=143\). It retains every correction
\(\chi_{-3}(p)p^{2e}a_{n/p}\) with \(p\mid n\), checks all coefficients
through \(n=37\) at \(p=7\) and \(n=20\) at \(p=13\), and constructs
the full integral matrices. The exact Faddeev--LeVerrier recursion
and three independent Bareiss determinant evaluations verify their
characteristic polynomials.

The unitary replay likewise reconstructs the full characteristic
polynomials, using exact fractions, and verifies the terminal
Cayley--Hamilton zero matrix. Their product decompositions have
degree patterns \(2+4+8\), \(2+5+14\), and \(2+5+20\), respectively
at both primes; only exact multiplication of those factors is used,
not an irreducibility claim. The complete expanded coefficient lists,
basis columns, and identifying matrices are in
\path{hecke_general/higher2/harmonic_hecke_exact_data.json};
the full classical outputs are in
\path{hecke_general/higher2/gl2_higher2_certificate.json}.

The separate standard-library comparison verifier evaluates those
complete polynomials by integer Horner arithmetic and computes the
two-value gcds. It verifies (HH.2) by exact multiplication.
Primality is proved recursively by the following elementary Lucas
criterion. For each integer \(n>2\), the certificate gives a complete
already-proved prime factorization of \(n-1\) and an integer \(a\)
such that
\[
 a^{n-1}\equiv1\pmod n,\qquad
 \gcd\bigl(a^{(n-1)/q}-1,n\bigr)=1
 \quad(q\mid n-1,\ q\text{ prime}).
\]
If \(r\) is any prime divisor of \(n\), these conditions force the
order of \(a\bmod r\) to be \(n-1\). Consequently \(n-1\mid r-1\),
so \(r\ge n\) and \(n=r\) is prime. Starting from \(2\), the finite
certificate verifies \(59\) such prime nodes, including every factor
in (HH.2). This is an exact primality proof, without reliance on a
probable-prime test or on a long trial-division search.
\end{proof}

\paragraph{The unitary factors alone and reproducibility.}
For comparison with the separate model reports, the unitary-only
common divisors are
\[
\begin{aligned}
\gcd_p P_{14,p}(E_{14,p})
 &=2^{70}3^{65}5^5\,53\,419\,15083\,16519\,7547263\,42800473,\\
\gcd_p P_{16,p}(E_{16,p})
 &=2^{81}3^{99}5^2\,11^2\,23\,401\,13687\,36293\,38177\\
 &\qquad{}\cdot25078104767246023,\\
\gcd_p P_{18,p}(E_{18,p})
 &=2^{125}3^{128}5^{15}\,13^2\,3343\,7691\,8609\,18541\\
 &\qquad{}\cdot151871\,8361790201,
\end{aligned}
\tag{HH.3}
\]
where each gcd is over \(p=7,13\).
The earlier model verifier is used unchanged, with its explicit
\texttt{--data}, \texttt{--neighbors}, and \texttt{--report} arguments
pointing to the new directory and the same frozen neighbor file.
All data, discovery wrappers, exact replay reports, the \(59\)-node
prime certificate, and the independent comparison verifier are under
\path{hecke_general/higher2/}.
The successful exact discovery times were \(2.78\), \(5.77\), and
\(14.44\) seconds for \(k=13,15,17\), respectively. The latter two
use modular screening of points after the initial symbolic
point-selection attempts were stopped.

This proposition supplies full finite models and exact integer
comparisons. An arithmetic deduction still uses the separate
admissibility, conductor, transfer, and trace interface, with an
auxiliary prime different from the residue characteristic. In
particular (HH.2) alone does not eliminate either of the two listed
large residual primes.
